%% Chapter 5 — Section 5.2: Full Conditionals and the Gibbs Sampler
%% A First Course in Monte Carlo Methods
%% D. Sanz-Alonso & O. Al-Ghattas

\section{Full Conditionals and the Gibbs Sampler}
\label{sec:5.2}

We next define the concept of full conditionals. For clarity, we give the definition in both
discrete and continuous cases.

\begin{definition}[Full Conditionals]
\label{def:5.1}
\begin{enumerate}[label=\roman*)]
  \item If $E$ is discrete and $f(x) = f(x_1, \ldots, x_d)$ is the p.m.f.\ of a random vector
    $X = (X_1, \ldots, X_d)$ taking values in $E \times \cdots \times E = E^d$. We denote
    \[
      f(x_{-j}) = \sum_{\xi \in E} f(x_1, \ldots, x_{j-1}, \xi, x_{j+1}, \ldots, x_d).
    \]
    The $j$-th full conditional p.m.f.\ is:
    \begin{align*}
      f_j(x_j | x_{-j})
        &= f_j(x_j | x_i, i \neq j)
         = \Prob(X_j = x_j \mid X_i = x_i, i \neq j)\\
        &= \frac{\Prob(X_j = x_j, X_i = x_i, i \neq j)}{\Prob(X_i = x_i, i \neq j)}\\
        &= \frac{f(x_1, \ldots, x_{j-1}, x_j, x_{j+1}, \ldots, x_d)}{f(x_{-j})}.
    \end{align*}

  \item If $E$ is continuous and $f(x) = f(x_1, \ldots, x_d)$ is the p.d.f.\ of the random
    vector $X = (X_1, \ldots, X_d)$ taking values on $E \times \cdots \times E = E^d$, we denote
    \[
      f(x_{-j}) = \int_E f(x_1, \ldots, x_{j-1}, \xi, x_{j-1}, \ldots, x_d)\, d\xi,
    \]
    and define the $j$-th full conditional density
    \[
      f_j(x_j | x_{-j}) = f_j(x_j | x_i, i \neq j)
      = \frac{f(x_1, \ldots, x_{j-1}, x_j, x_{j+1}, \ldots, x_d)}{f(x_{-j})}.\qquad\square
    \]
\end{enumerate}
\end{definition}

The distribution defined by $f_j(x_j | x_i, i \neq j)$ is called the $j$-th full conditional
distribution of $f$. We will only need to know these distributions up to a normalizing constant,
and so terms that do not depend on $x_j$ can be ignored. In particular, it is useful to note that
full conditionals are proportional to the joint density.

\begin{example}[Full Conditionals of Multivariate Gaussians]
\label{ex:5.2}
Consider a multivariate Gaussian random vector $X = (X_1, \ldots, X_d) \sim \Normal(\mu, H^{-1})$.
Then, $X_j | X_{-j} \sim \Normal(\nu_j, H_{jj}^{-1})$, with
\[
  \nu_j = \mu_j - H_{jj}^{-1} \sum_{k \neq j} H_{jk}(x_k - \mu_k).
\]
To check this, recall that a univariate normal random variable with mean $\alpha$ and precision
$\beta$ has p.d.f.\ proportional to
\begin{equation}
  \exp\!\Bigl(-\tfrac{1}{2}x^2\beta + x\beta\alpha\Bigr).
  \label{eq:5.1}
\end{equation}
Assume for the moment that $\mu = 0 \in \R^d$. Using that the $j$-th full conditional is
proportional to the joint density,
\begin{equation}
  f_j(x_j | x_{-j})
  \propto \exp\!\Bigl(-\tfrac{1}{2}x^T H x\Bigr)
  \propto \exp\!\Bigl(-\tfrac{1}{2}x_j^2 H_{jj} - x_j \sum_{k \neq j} H_{jk} x_k\Bigr).
  \label{eq:5.2}
\end{equation}
Thus $X_j | X_{-j}$ is univariate normal. Comparing the coefficients of the quadratic terms in
Equations~\eqref{eq:5.1} and~\eqref{eq:5.2} we see that the precision is $H_{jj}$, and
comparing the coefficients for the linear term we see that the mean is $\nu_j$. Finally, if $X$
has mean $\mu$, then applying the above argument to $X - \mu$ gives the result.
\end{example}

The Gibbs sampler is based on sampling full conditionals. It is not immediately obvious that full
conditionals (unlike marginals) can specify a distribution uniquely, provided they are consistent.
Sufficient and necessary conditions are given by the Hammersley--Clifford theorem, which for
simplicity we only present in the 2-dimensional case.

\begin{theorem}[Hammersley--Clifford 2-dimensional Case]
\label{thm:5.3}
If $\int \frac{f(x_2|x_1)}{f(x_1|x_2)}\,dx_2$ exists, then the joint density associated with
$f(x_2|x_1)$ and $f(x_1|x_2)$ is given by
\[
  f(x_1, x_2) = \frac{f(x_2|x_1)}{\displaystyle\int \frac{f(x_2|x_1)}{f(x_1|x_2)}\,dx_2}.
\]
\end{theorem}

\begin{proof}
Since $f(x_2|x_1)f(x_1) = f(x_1|x_2)f(x_2)$,
\[
  \int \frac{f(x_2|x_1)}{f(x_1|x_2)}\,dx_2
  = \int \frac{f(x_2)}{f(x_1)}\,dx_2
  = \frac{1}{f(x_1)},
\]
and so the claim follows if the integral above exists.
\end{proof}

We are now ready to introduce the Gibbs sampler.

\begin{algorithm}[H]
\caption{Gibbs Sampler}
\label{alg:5.1}
\begin{algorithmic}[1]
\Require target $f(x_1, \ldots, x_d)$, initialization
  $X^{(0)} = (X_1^{(0)}, \ldots, X_d^{(0)})$, sample size $N$.
\For{$n = 1, \ldots, N$}
  \State Choose $j \in \{1, \ldots, d\}$ (randomly or sequentially).
  \State Sample $X_j^{(n)} \sim f_j(x_j | X_i^{(n-1)}, i \neq j)$ and set
    \[
      X^{(n)} = \Bigl(X_1^{(n-1)}, \ldots, X_{j-1}^{(n-1)},\,
        X_j^{(n)},\, X_{j+1}^{(n-1)}, \ldots, X_d^{(n-1)}\Bigr).
    \]
\EndFor
\Ensure Sample $\{X^{(n)}\}_{n=1}^N$.
\end{algorithmic}
\end{algorithm}

\medskip
\noindent\fbox{\parbox{\dimexpr\linewidth-2\fboxsep-2\fboxrule\relax}{%
\textbf{Pros and Cons}: An important advantage of the Gibbs sampler is that there are no
parameters to choose and tune. The Gibbs sampler is particularly useful in problems where the
target distribution cannot be sampled directly, but its full conditionals are easy to sample
from; e.g.\ in hierarchical Bayesian models, mixture models, and Bayesian missing data problems.
Updating only one variable at a time makes the implementation scalable to high-dimensional
settings. However, the convergence can be slow when the individual variables are strongly
correlated. Moreover, as other methods based on Markov chain sampling, the Gibbs sampler is not
parallelizable.
}}

\medskip

There are two main ways to improve the performance of the Gibbs sampler when the variables are
correlated.

\begin{enumerate}
  \item Reparametrize the distribution so that the new variables are less correlated.

  \item Block variables that are highly correlated. Suppose for example that $X_1, X_2$ are
    highly correlated and we want to sample from the distribution of $(X_1, X_2, X_3)$. Then,
    we could do:
    \begin{itemize}
      \item Sample $X_1, X_2 | X_3$.
      \item Sample $X_3 | X_1, X_2$.
    \end{itemize}
\end{enumerate}

We conclude this section showing that the Gibbs sampler is a Metropolis Hastings algorithm.

\begin{theorem}[Gibbs Sampler is a Metropolis Hastings Algorithm]
\label{thm:5.4}
The Gibbs sampler is a Metropolis Hastings algorithm that uses full conditionals as proposals.
Proposed moves are always accepted.
\end{theorem}

\begin{proof}
Let $x = (x_1, \ldots, x_d)$ and $z = (z_1, \ldots, z_d)$ with $x_i = z_i$ for $i \neq j$.
\begin{align*}
  a(x, z)
  &= \min\!\left\{1,\, \frac{f(z)}{f(x)}\frac{f_j(z_j | z_i, i \neq j)}{f_j(x_j | x_i, i \neq j)}\right\}\\
  &= \min\!\left\{1,\, \frac{f(z)}{f(x)}\frac{f(z)/f(z_{-j})}{f(x)/f(x_{-j})}\right\}
    \qquad\text{(by definition of full conditional)}\\
  &= \min\!\left\{1,\, \frac{f(z)}{f(x)}\frac{f(x)/f(x_{-j})}{f(z)/f(x_{-j})}\right\}
    \qquad(f(z_{-j}) = f(x_{-j}) \text{ because } x_i = z_i\ \forall i \neq j)\\
  &= 1.
\end{align*}
\end{proof}

The Gibbs sampler is a particular example of a Metropolis Hastings algorithm, and therefore its
transition kernel satisfies detailed balance with respect to $f$. We show this directly as an
exercise in the following proposition.

\begin{proposition}
\label{prop:5.5}
Each transition of the Gibbs sampler satisfies detailed balance with respect to the target
distribution $f$.
\end{proposition}

\begin{proof}
Denote by $p_{\mathrm{GS}}$ the transition kernel of the Gibbs sampler and let
\begin{align*}
  x &= (x_1, \ldots, x_d),\\
  z &= (x_1, \ldots, x_{j-1}, \xi, x_{j-1}, \ldots, x_d).
\end{align*}
Then,
\begin{align*}
  f(x)\,p_{\mathrm{GS}}(x, z)
  &= f(x)\,\frac{f(z)}{f(z_{-j})}
   = f(z)\,\frac{f(x)}{f(x_{-j})}\\
  &= f(z)\,f_j(x_j | x_i, i \neq j)
   = f(z)\,p_{\mathrm{GS}}(z, x).
\end{align*}
\end{proof}

\begin{example}[2D Slice Sampler]
\label{ex:5.6}
Recall that when we introduced rejection sampling in Chapter~2 we observed that given a density
$f(x)$, $x \in \R$, we can represent it as the marginal of $f(x, u) = \mathbf{1}_{\{0 < u < f(x)\}}$.

\begin{figure}[htbp]
  \centering
  \includegraphics[width=0.75\textwidth]{figures/fig_05_1}
  \caption{\textit{$f(x, u) = \mathbf{1}_{\{0 < u < f(x)\}}$.}}
  \label{fig:5.1}
\end{figure}

The full conditionals are
\begin{align*}
  f(x|u) &\propto \mathbf{1}_{\{0 < u < f(x)\}},\\
  f(u|x) &\propto \mathbf{1}_{\{0 < u < f(x)\}}.
\end{align*}
We can use the Gibbs sampler to obtain samples from $f(x, u)$. Keeping the $x$-variable, we
obtain samples from $f(x)$. This specific type of Gibbs sampler is called the 2D slice sampler.
\end{example}
